\section{Introduction}

Active mobility — travel on foot, by bicycle, or by any means other than individually motorized vehicles — can help mitigate climate impacts. According to \cite{BRAND2021102764}, about 70\% of the life-cycle CO2 emissions per person from urban mobility activities come from cars, against only 1\% from bicycles. The authors estimate that each avoided car trip reduces individual emissions by 62\%.

Urban planning models have emerged in response to the climate emergency and to the search for decarbonization policies. One prominent model, currently being implemented in cities such as Paris and Barcelona, is the 15-minute city \cite{Allam2022-yx}. It organizes the city so that the services residents use daily are reachable through short trips, prioritizing active mobility.

The viability of such models depends on how the travel network behaves under real operating conditions. Prior work has shown that urban accessibility is not uniform: socioeconomic inequalities strongly shape who reaches which services and in what time \cite{sbbd_estendido}. Extreme climate events such as floods, intense rainfall, and landslides degrade mobility infrastructure, interrupting or obstructing segments of the road network and rendering certain trips unfeasible or substantially more costly. These disruptions do not affect the population evenly. Areas with higher vulnerability often show lower resilience to network degradation, and populations with less adaptive capacity bear larger impacts.

Answering these questions requires computing accessibility over heterogeneous spatial data. A city is described by a street graph, a set of points of interest, a discrete territorial grid, environmental rasters or polygons, and census aggregates, each with its own geometry, resolution, and coordinate reference system. Accessibility is then a many-to-many shortest-path problem whose cost function must absorb environmental evidence that is continuous in some cities and categorical in others. The computational problem is to make that pipeline reproducible and city-agnostic, so that the same analysis transfers to a new municipality by changing configuration rather than code.

This work evaluates how climate disruptions change urban accessibility and how those impacts are distributed unevenly across space and socioeconomic profiles. We propose a methodological approach that integrates travel-network modeling with socioeconomic analysis and compares typical scenarios (without disruption) against scenarios conditioned by environmental variables. The approach is implemented as an open-source library with a declarative workflow, in which penalty functions, data layers, and scenario definitions are configured rather than hard-coded.

The contributions of this work are threefold:

\begin{itemize}
  \item an open computational framework for modeling urban accessibility under environmental disruption, whose region-agnostic data model and declarative scenario configuration allow the same pipeline to be applied to different cities;
  \item a comparison of scenarios that distinguishes access loss from travel time increase, producing quantitative evidence on two Brazilian state capitals under different climate stressors;
  \item an assessment of how accessibility degradation correlates with socioeconomic conditions, exposing which populations bear the largest impacts.
\end{itemize}

The remainder of this paper is organized as follows. Section~\ref{sec:fundamentals} reviews the concepts of urban accessibility, active mobility, and environmental disruption. Section~\ref{sec:methodology} presents the framework and its penalization mechanism, and Section~\ref{sec:case_studies} describes the two case studies and their calibration. Section~\ref{sec:results} reports the results. Section~\ref{sec:discussion} discusses limitations and implications, and Section~\ref{sec:conclusion} concludes.
